无线电干扰源的定位与清除是频谱管理的重要任务，其效率取决于检测点选取与行进路线安排。本文针对此问题，建立\textbf{半平面交定位模型}、\textbf{基于最坏定位误差的第二检测点选址模型}与\textbf{搜索-清除联合规划模型}，给出定位区域计算、检测点选择与全区域清除的策略与算法。

\textbf{针对问题一：}我们建立了\textbf{半平面交定位模型}。首先，依据示向度及其±1°误差界，将每次观测转化为两个半平面约束，联立\textbf{各次观测的半平面约束}得到角度交会区域即定位区域。接着，构建\textbf{半平面求交与顶点枚举算法}计算区域直径。随后，通过构造等边三角形，求得其最小包围圆半径\textbf{大于区域直径的一半}，说明以区域直径为直径的圆\textbf{不一定覆盖定位区域}。最后，构造多种实例验证算法，程序返回区域均与预设结果一致，体现了算法的\textbf{正确性与鲁棒性}。

\textbf{针对问题二：}我们建立了\textbf{基于最坏定位误差的第二检测点选址模型}。首先，在第一问的基础上考虑\textbf{有效接收距离}，推导三类反馈下的定位区域更新规则，将\textbf{能获得信息增益的点集作为候选区域}。接着，采用\textbf{最小包围圆估计}方法，以圆心预测干扰源位置、以半径衡量最坏预测误差；依据\textbf{极小化极大准则}，对每个候选检测位置，取其在最不利反馈下的预测误差；再选取该误差最小的位置作为最优二次检测点。求解时采用\textbf{粗网格与局部增密搜索算法}，得到数值最优检测点。最后，在大量模拟场景上开展\textbf{配对实验}，结果显示：预测位置与真值的均方根误差为\textbf{12.091米}，低于随机选点的229.057米，表明其在该测试分布下能够\textbf{有效提高定位精度}。

\textbf{针对问题三：}我们建立了\textbf{基于覆盖约束的搜索-清除联合规划模型}。借鉴\textbf{马尔可夫决策过程}的序贯决策思想，以机器狗位置、当前频道、已发现频道的位置—接收半径组合集与未发现频道的未排除区域为状态变量，建立\textbf{状态转移关系与双终止条件}，并以\textbf{完整任务时间的期望最小}为目标，在全部清除干扰源的情况下，按检测与清除反馈动态更新状态与优化行动策略。求解时，借鉴\textbf{模型预测控制}的滚动优化思想，采用了\textbf{基于条件场景评估的启发式滚动规划算法}，通过\textbf{多起点局部搜索、两交换与任务重插入}改进行动次序，每次仅执行首个行动，并根据反馈重新规划。在1000个独立模拟案例中，\textbf{平均每源耗时为212.77秒}。此外，三次正式测试\textbf{平均每源完整退出耗时为213.81秒}。

\textbf{针对问题四：} 我们建立了\textbf{方向覆盖约束下的搜索-清除联合规划模型}。针对发射朝向未知导致单点无信号无法直接排除区域的问题，将第三问的单点圆盘覆盖排除条件扩展为\textbf{距离与方向联合排查条件}，利用接收距离限制与检测点凸包包围关系，建立与源朝向无关的\textbf{区域排除依据}，并以\textbf{完整任务时间的期望最小}为目标，在保证完整排查的约束下优化行动策略。求解算法采用\textbf{基于有限参数组合估计的启发式滚动规划算法}，预测后续检测次数与定位代价，通过\textbf{迭代局部搜索与检测点位置调整}协同优化巡测布局和任务顺序，并根据实际反馈更新规划。结果显示：\textbf{平均每源完整退出耗时为416.13秒}。此外，三次正式测试\textbf{平均每源完整退出耗时为458.77秒}。

本文创新性在于：一是构建了有界示向度约束下的鲁棒定位与检测点选择模型；二是结合了序贯决策与启发式滚动规划，实现搜索、定位与清除的一体化决策；三是针对了定向源的无信号歧义，建立多测站距离—方向联合排查方法，并将其融入在线协同规划。
